\subsection{欧拉–麦克劳林求和公式余项的估计}

在 (16) 中对区间 \([0, 1]\) 上的伯努利多项式得到的估计式，使我们能容易地估计欧拉–麦克劳林求和公式（VI, p.~282，公式 (39)）的余项 \(T_{p}(x, n)\) ：

\[
\left\{ \begin{array}{l} f (x) + f (x + 1) + \dots + f (x + n) \\ = \int_ {x} ^ {x + n + 1} f (t) d t - \frac {1}{2} \bigl (f (x + n + 1) - f (x) \bigr) \\ \qquad + \sum_ {k = 1} ^ {p} \frac {b _ {2 k}}{(2 k) !} \bigl (f ^ {(2 k - 1)} (x + n + 1) - f ^ {(2 k - 1)} (x) \bigr) + T _ {p} (x, n). \end{array} \right.\tag{1}
\]

事实上，有（VI, p.~283，公式 (41)）

\[
\mathrm{T} _ {p} (x, n) = - \frac {1}{(2 p + 1) !} \int_ {0} ^ {n + 1} \overline {{\mathrm{B}}} _ {2 p + 1} (1 - s) f ^ {(2 p + 1)} (x + s) d s\tag{2}
\]

其中 \(\overline{\mathbf{B}}_{2p + 1}(t)\) 是以 1 为周期、等于 \(\mathrm{B}_{2p + 1}(t)\) （在区间 {[}0, 1{[} 上）的周期函数。VI, p.~288 的公式 (16) 表明

\[
\left| \overline {{{{\mathrm{B}}}}} _ {2 p + 1} (t) \right| \leqslant 4 e ^ {2 \pi} \frac {(2 p + 1) !}{(2 \pi) ^ {2 p + 1}}\tag{3}
\]

对一切 \(t \in R\) 成立，而应用中值公式即得估计式

\[
\left| \mathrm{T} _ {p} (x, n) \right| \leqslant \frac {4 e ^ {2 \pi}}{(2 \pi) ^ {2 p + 1}} \int_ {x} ^ {x + n + 1} \left| f ^ {(2 p + 1)} (t) \right| d t\tag{4}
\]

对 \(\mathrm{T}_p(x,n)\) 成立。

